\documentclass[12pt]{article}
\usepackage[spanish,es-nodecimaldot]{babel}
\usepackage[utf8]{inputenc}
\usepackage[T1]{fontenc}
\usepackage[a4paper,margin=2.5cm]{geometry}
\usepackage{lmodern}
\usepackage{setspace}
\usepackage{parskip}
\usepackage{graphicx}
\usepackage{float}
\usepackage{amsmath}
\setlength{\belowcaptionskip}{6pt}

\begin{document}
\thispagestyle{empty}

\begin{center}
    \textsc{Universidad Central de Venezuela}\\[0.35cm]
    \textsc{Facultad de Ingenier\'ia}\\[0.35cm]
    \textsc{Escuela de Ingenier\'ia El\'ectrica}\\[0.35cm]
    \textsc{Departamento de Potencia}\\[0.35cm]
    \textsc{Laboratorio de M\'aquinas El\'ectricas I}
\end{center}

\vspace{6.8cm}

\begin{center}
    \textbf{PRE-LABORATORIO N\textsuperscript{o} 8}\\[0.35cm]
    \textbf{Transformador monof\'asico en Paralelo}
\end{center}

\vfill

\noindent
\begin{minipage}[t]{0.48\textwidth}
    Docente:\\[0.25cm]
    Jes\'us Blondel
\end{minipage}%
\hfill
\begin{minipage}[t]{0.48\textwidth}
    \raggedleft
    Autor:\\[0.25cm]
    Br. Carla Fajardo\\[0.25cm]
    V-27.571.576
\end{minipage}

\vspace{1.1cm}

\begin{center}
    Caracas, \today
\end{center}

\newpage

\section{Introducci\'on}
Este prelaboratorio estudia el comportamiento de un banco trif\'asico de transformadores monof\'asicos en conexi\'on estrella--estrella (Y--Y). Se analiza la influencia del neutro y del devanado terciario en delta sobre las corrientes de excitaci\'on, los arm\'onicos y las tensiones del sistema, tanto en vac\'io como ante cargas lineales y no lineales, por lo que se comparan las condiciones de neutro conectado y aislado, junto con la funci\'on del devanado terciario como trayectoria para las corrientes arm\'onicas. El registro de tensiones de fase y de l\'inea, y de corrientes de excitaci\'on y del neutro, permite relacionar las formas de onda observadas con el tipo de carga y con la alimentaci\'on a tres o cuatro hilos. As\'i, se busca interpretar las distorsiones y variaciones de tensi\'on que aparecen en cada condici\'on de operaci\'on.
\newpage

\section{Objetivos}

\subsection*{Objetivo general}
Estudiar el comportamiento de un banco de transformaci\'on trif\'asico conformado por unidades monof\'asicas, en conexi\'on Estrella-Estrella (Y--Y), analizando la influencia de la conexi\'on de neutro y del delta terciario sobre las formas de onda de tensi\'on y corriente, tanto en vac\'io como con cargas lineales y no lineales.

\subsection*{Objetivos espec\'ificos}
\begin{enumerate}
    \item Observar las formas de onda de tensi\'on y corriente de un banco de transformaci\'on trif\'asico en conexi\'on Estrella-Estrella (Y--Y).
    \item Comparar el comportamiento del banco Y--Y alimentado con un sistema de cuatro hilos y con un sistema de tres hilos.
    \item Registrar y analizar las formas de onda de tensi\'on fase-neutro y fase-fase, as\'i como la corriente de excitaci\'on, en condiciones de vac\'io con y sin neutro conectado.
    \item Analizar la caracter\'istica de tensi\'on terminal del secundario en funci\'on de la corriente de carga, $U=f(I)$, para distintas condiciones de carga:
    \begin{itemize}
        \item Carga trif\'asica lineal balanceada con sistema de cuatro hilos.
        \item Carga trif\'asica lineal balanceada con sistema de tres hilos.
    \end{itemize}
    \item Analizar el comportamiento del banco trif\'asico alimentando cargas no lineales:
    \begin{itemize}
        \item Carga trif\'asica no lineal balanceada mediante rectificador trif\'asico.
        \item Carga monof\'asica no lineal mediante rectificador monof\'asico de onda completa.
    \end{itemize}
\end{enumerate}

\newpage

\section{Marco te\'orico}
\subsection{Conexiones de transformadores trif\'asicos: generalidades}

La transformaci\'on de tensiones y corrientes en los sistemas trif\'asicos puede llevarse a cabo mediante dos arquitecturas constructivas:
\begin{itemize}
    \item \textbf{Banco trif\'asico formado por tres transformadores monof\'asicos independientes:} ofrece mayor facilidad de transporte y mantenimiento (un transformador de reserva puede respaldar al banco).
    \item \textbf{Transformador trif\'asico de n\'ucleo com\'un:} incorpora los tres conjuntos de devanados sobre una misma estructura magn\'etica (generalmente de tres columnas). Es m\'as compacto, ligero, econ\'omico y ligeramente m\'as eficiente.
\end{itemize}

Los devanados primarios y secundarios de estos equipos pueden interconectarse de forma independiente en diversas configuraciones:
\begin{itemize}
    \item \textbf{Estrella--Estrella ($Y$--$Y$ o $Yy$):} ambos lados conectados en estrella.
    \item \textbf{Estrella--Delta ($Y$--$\Delta$ o $Yd$):} primario en estrella y secundario en tri\'angulo o delta.
    \item \textbf{Delta--Estrella ($\Delta$--$Y$ o $Dy$):} primario en tri\'angulo y secundario en estrella.
    \item \textbf{Delta--Delta ($\Delta$--$\Delta$ o $Dd$):} ambos lados conectados en tri\'angulo.
    \item \textbf{Zig-zag ($Y$--$z$ o $Yz$):} subdivide las bobinas secundarias en dos partes iguales conectadas en serie-oposici\'on con fases adyacentes para equilibrar cargas monof\'asicas.
\end{itemize}

Existen tambi\'en esquemas especiales con dos transformadores, como la conexi\'on delta abierta ($V$--$V$) y la conexi\'on Scott--T.

\subsection{Conexi\'on Estrella--Estrella ($Y$--$Y$)}

En una conexi\'on Estrella--Estrella, las relaciones entre las magnitudes de fase y de l\'inea son:
\begin{itemize}
    \item Tensi\'on de l\'inea: $V_L=\sqrt{3}\,V_f$.
    \item Corriente de l\'inea: $I_L=I_f$.
\end{itemize}

\subsubsection*{Ventajas}
\begin{itemize}
    \item \textbf{Menor exigencia de aislamiento:} la tensi\'on aplicada a cada bobina es $V_L/\sqrt{3}$ (un 58\% de la tensi\'on entre l\'ineas). Esto reduce los esfuerzos diel\'ectricos sobre el aislamiento, haci\'endola muy ventajosa en alta y muy alta tensi\'on.
    \item \textbf{Mayor robustez mec\'anica:} requiere conductores de mayor secci\'on transversal frente a un devanado en delta equivalente, resultando en bobinas mec\'anicamente m\'as resistentes a los esfuerzos de cortocircuito.
    \item \textbf{Sin desplazamiento angular:} no introduce desfase entre las tensiones del primario y del secundario ($0^\circ$ o $180^\circ$, \'indice horario 0 o 6).
\end{itemize}

\subsubsection*{Inconvenientes principales}
\begin{enumerate}
    \item \textbf{Sensibilidad a cargas desequilibradas:} si las cargas conectadas entre fase y neutro no est\'an balanceadas, se generan desplazamientos en el potencial del neutro.
    \item \textbf{Generaci\'on de tensiones arm\'onicas de tercer orden:} presencia de componentes significativos del tercer arm\'onico por la saturaci\'on del n\'ucleo magn\'etico.
\end{enumerate}

\subsection{Influencia del neutro en la conexi\'on Estrella--Estrella}

El comportamiento del transformador $Y$--$Y$ cambia de manera radical seg\'un se disponga o no de la conexi\'on del neutro.

\subsubsection*{A) Con neutro aislado o sin conductor neutro (sistema a 3 hilos)}
\begin{itemize}
    \item \textbf{Condici\'on el\'ectrica:} al no existir un camino de retorno, la suma instant\'anea de las tres corrientes de excitaci\'on debe ser nula en el punto neutro: $i_a+i_b+i_c=0$.
    \item \textbf{Comportamiento arm\'onico:} como los terceros arm\'onicos de la corriente de excitaci\'on est\'an en fase en las tres ramas (componentes homopolares), no pueden circular. Esto obliga a que la corriente de excitaci\'on sea puramente senoidal.
    \item \textbf{Deformaci\'on del flujo:} al ser la corriente de excitaci\'on senoidal en un n\'ucleo saturado, el flujo magn\'etico se deforma y pasa a contener una fuerte componente de tercer arm\'onico ($\Phi_3$).
    \item \textbf{Efecto sobre las tensiones y oscilaci\'on del neutro:} este flujo distorsionado induce fuerzas electromotrices con un gran contenido de tercer arm\'onico ($E_3$) en cada fase.
    \begin{itemize}
        \item En las tensiones compuestas (de l\'inea a l\'inea), el tercer arm\'onico se anula: $E_{a3}-E_{b3}=0$.
        \item En las tensiones simples (fase a neutro), los terceros arm\'onicos de cada fase se suman al fundamental. Esto produce sobretensiones en las fases y hace que el punto neutro del secundario oscile en una trayectoria circular a la frecuencia del tercer arm\'onico, fen\'omeno conocido como ``neutro inquieto'' u oscilaci\'on del neutro.
    \end{itemize}
\end{itemize}

\subsubsection*{B) Con neutro primario r\'igidamente conectado a tierra (sistema a 4 hilos)}
\begin{itemize}
    \item El conductor neutro proporciona un camino de retorno para la suma de las componentes de tercer arm\'onico de la corriente de excitaci\'on ($i_{N,3}=3i_{o,3}$).
    \item Al permitir que la corriente de excitaci\'on contenga la componente acampanada de tercer arm\'onico que requiere el n\'ucleo, el flujo magn\'etico se mantiene sinusoidal y las tensiones de fase permanecen equilibradas y sin distorsi\'on.
    \item \textit{Desventaja:} la circulaci\'on de corrientes arm\'onicas por el hilo neutro y las l\'ineas a\'ereas puede provocar ca\'idas de tensi\'on distorsionantes e interferencias inductivas en l\'ineas de comunicaci\'on o telef\'onicas cercanas.
\end{itemize}

\subsection{Devanado delta terciario (terciario en $\Delta$)}

Para resolver los problemas de distorsi\'on y desequilibrio en la conexi\'on $Y$--$Y$ sin necesidad de conectar el neutro primario a la red de alimentaci\'on, se a\~nade un tercer devanado independiente conectado en tri\'angulo ($\Delta$).

\subsubsection*{Principio de funcionamiento}
\begin{enumerate}
    \item Las componentes de fuerza electromotriz de tercer arm\'onico ($E_3$) inducidas en cada una de las tres fases del devanado terciario est\'an en fase.
    \item En la conexi\'on en tri\'angulo, estas fuerzas electromotrices no quedan bloqueadas, sino que se suman y producen una fuerza electromotriz resultante nula hacia el exterior, pero con un valor neto de $3E_3$ dentro del lazo delta.
    \item Esta fuerza electromotriz hace circular una corriente de tercer arm\'onico atrapada dentro del delta.
    \item Dicha corriente crea una fuerza magnetomotriz en oposici\'on que anula o ``borra'' el flujo de tercer arm\'onico en el n\'ucleo ($\Phi_3\approx 0$).
\end{enumerate}

\subsubsection*{Beneficios}
\begin{itemize}
    \item Restablece una forma de onda senoidal en los flujos del n\'ucleo y en las tensiones inducidas.
    \item Proporciona un camino de baja impedancia para equilibrar cargas desbalanceadas en el secundario, atenuando la traslaci\'on del neutro.
    \item \textbf{Capacidad y aplicaciones adicionales:} por lo general, se dise\~na para soportar alrededor de un tercio (33\%) de la potencia nominal del transformador. A menudo se sacan sus bornes al exterior para alimentar servicios auxiliares de la subestaci\'on o para conectar capacitores o condensadores s\'incronos de regulaci\'on de tensi\'on y factor de potencia.
\end{itemize}

\subsection{Cargas lineales y no lineales}
\begin{itemize}
    \item \textbf{Cargas lineales:} elementos en los que la impedancia es constante e independiente de la tensi\'on aplicada (resistencias puras, inductancias y capacitancias). Si la tensi\'on de alimentaci\'on es senoidal, la corriente absorbida ser\'a de la misma frecuencia y con una forma de onda puramente senoidal.
    \item \textbf{Cargas no lineales:} dispositivos cuya impedancia var\'ia durante el ciclo el\'ectrico o que presentan comportamientos no lineales (convertidores est\'aticos, rectificadores de diodos o tiristores, variadores de frecuencia, balastos y transformadores operados en la zona de saturaci\'on).
    \begin{itemize}
        \item Estas cargas absorben impulsos de corriente distorsionados.
        \item Mediante el desarrollo en series de Fourier, estas corrientes se descomponen en una componente fundamental m\'as un espectro de corrientes arm\'onicas que se inyectan hacia la red el\'ectrica.
    \end{itemize}
\end{itemize}

\subsection{Corriente de excitaci\'on y arm\'onicos}

La corriente que absorbe un transformador en vac\'io se denomina corriente de excitaci\'on ($i_{\mathrm{exc}}$), y consta de dos componentes principales:
\begin{enumerate}
    \item \textbf{Corriente de p\'erdidas en el n\'ucleo ($i_{h+e}$):} est\'a en fase con la tensi\'on y cubre las p\'erdidas por hist\'eresis y por corrientes par\'asitas o de Foucault.
    \item \textbf{Corriente magnetizante ($i_m$):} est\'a desfasada $90^\circ$ con respecto a la tensi\'on y se encarga de establecer el flujo magn\'etico en el hierro.
\end{enumerate}

\subsubsection{Origen de los arm\'onicos por la curva de imanaci\'on}

Debido a la naturaleza no lineal y al fen\'omeno de saturaci\'on de la curva $B$--$H$ de los materiales ferromagn\'eticos, para mantener un flujo magn\'etico senoidal en el n\'ucleo, la corriente magnetizante necesaria no puede ser senoidal, sino que adopta una forma acampanada con picos pronunciados.

Al descomponer esta onda acampanada mediante series de Fourier, aparecen arm\'onicos impares ($1.\textsuperscript{er}$, $3.\textsuperscript{er}$, $5.\textsuperscript{o}$, $7.\textsuperscript{o}$, \ldots):
\begin{itemize}
    \item \textbf{Primer arm\'onico (fundamental):} frecuencia industrial (50 Hz o 60 Hz).
    \item \textbf{Tercer arm\'onico:} es la componente arm\'onica m\'as prominente y representa habitualmente entre el 40\% y el 45\% del total de la corriente de excitaci\'on en transformadores t\'ipicos.
    \item \textbf{Quinto y s\'eptimo arm\'onicos:} representan t\'ipicamente alrededor del 15\% y el 3\%, respectivamente.
\end{itemize}

\subsubsection{Distribuci\'on de arm\'onicos en sistemas trif\'asicos}
\begin{itemize}
    \item \textbf{Arm\'onicos triplen ($3.\textsuperscript{o}$, $9.\textsuperscript{o}$, $15.\textsuperscript{o}$, \ldots):} est\'an en fase en las tres ramas de un sistema trif\'asico equilibrado (homopolares). En conexiones $\Delta$ circulan en el bucle cerrado sin salir a las l\'ineas; en conexiones $Y$ sin neutro quedan bloqueados en corriente, distorsionando la forma del flujo.
    \item \textbf{Arm\'onicos del grupo ($5.\textsuperscript{o}$, $7.\textsuperscript{o}$, $11.\textsuperscript{o}$, $13.\textsuperscript{o}$, \ldots):} tienen secuencias de fase positiva y negativa, transmiti\'endose libremente entre las l\'ineas del sistema trif\'asico.
\end{itemize}

\newpage
\section{Equipos e instrumentos}
\begin{itemize}
    \item 3 transformadores monof\'asicos.
    \item Variac
    \item Osciloscopio
    \item Volt\'imetro de hierro m\'ovil.
    \item Amper\'imetro de hierro m\'ovil.
    \item Re\'ostato.
    \item Diodos.
    \item Resistencia shunt ($R_s$).
    \item Fuente de alimentaci\'on AC.
    \item Carga lineal.
    \item Carga no lineal.
    \item Cables de conexi\'on.
    \item Interruptores de prueba.
\end{itemize}

\newpage

\section{Diagramas de conexi\'on}

\begin{figure}[H]
    \centering
    \includegraphics[width=0.9\textwidth]{\detokenize{imagenes/conexion Y-Y.JPG}}
    \caption{Conexi\'on Y--Y.}
    \label{fig:conexion-yy}
\end{figure}

\begin{figure}[H]
    \centering
    \includegraphics[width=0.5\textwidth]{imagenes/carga\space no\space lineal\space monofásica.JPG}
    \caption{Carga no lineal monof\'asica.}
    \label{fig:carga-no-lineal-monofasica}
\end{figure}

\begin{figure}[H]
    \centering
    \includegraphics[width=0.5\textwidth]{imagenes/carga\space no\space lineal\space trifásica.JPG}
    \caption{Carga no lineal trif\'asica.}
    \label{fig:carga-no-lineal-trifasica}
\end{figure}

\section{Metodolog\'ia}

Antes de dar inicio a las mediciones, es imprescindible recopilar los datos t\'ecnicos y de placa de cada una de las unidades monof\'asicas que conformar\'an el banco trif\'asico. Esta informaci\'on resulta determinante para dise\~nar los diagramas de conexi\'on, definir la secuencia de operaciones y fijar las condiciones bajo las cuales se realizar\'an las pruebas.

Toda conexi\'on del circuito a la red de alimentaci\'on el\'ectrica deber\'a efectuarse a trav\'es de un dispositivo de protecci\'on (breaker). La energizaci\'on de cualquier arreglo experimental quedar\'a supeditada a la autorizaci\'on previa del profesor responsable.

\subsection{Parte: Importancia de la conexi\'on del neutro}

En este apartado se examinar\'a la conexi\'on Y--Y y se tomar\'an las formas de onda correspondientes a la corriente de excitaci\'on y a las tensiones fase-neutro y fase-fase. 

\begin{enumerate}
    \item \textbf{Condici\'on de vac\'io:}
    \begin{itemize}
        \item Se obtendr\'an las formas de onda de las tensiones y corrientes de inter\'es con el banco trif\'asico conectado en vac\'io y con el neutro interconectado entre la fuente y el primario.
        \item Posteriormente, se repetir\'an el paso anterior pero suprimiendo la conexi\'on del neutro entre la fuente y el primario.
        \item Se analizar\'an las formas de onda obtenidas y se explicar\'a la importancia de la conexi\'on del neutro en el primario.
    \end{itemize}
    \item \textbf{Con carga lineal trif\'asica balanceada:}
    \begin{itemize}
        \item Se determinar\'a la curva de tensi\'on terminal en el secundario en funci\'on de la corriente de carga, $U=f(I)$, para la conexi\'on Y--Y con carga trif\'asica lineal balanceada y alimentaci\'on a cuatro hilos.
        \item Se determinar\'a igualmente la curva $U=f(I)$ para la conexi\'on Y--Y con carga trif\'asica lineal balanceada y alimentaci\'on a tres hilos. Adem\'as, en condici\'on de plena carga, se registrar\'an las formas de onda de las tensiones y corrientes de inter\'es.
    \end{itemize}
    \item \textbf{Con carga monof\'asica lineal:}
    \begin{itemize}
        \item Para las conexiones Y--Y con carga monof\'asica lineal, se obtendr\'a la curva de tensi\'on terminal del secundario en funci\'on de la corriente de carga, $U=f(I)$, en la fase que soportar\'a la carga, as\'i como las tensiones presentes en las fases restantes, considerando que:
        \begin{itemize}
            \item Con neutro conectado en el primario.
            \item Sin neutro en el primario.
            \item Con el delta terciario conectado.
        \end{itemize}
        \item En condici\'on de plena carga, se registrar\'an las formas de onda de las tensiones y corrientes de inter\'es.
    \end{itemize}
\end{enumerate}

\subsection{Parte II: Corriente en el neutro y tensi\'on entre neutros}

\begin{enumerate}
    \item \textbf{Conexi\'on Y--Y alimentando una carga trif\'asica no lineal:}
    \begin{itemize}
        \item Se capturar\'a la forma de onda de las corrientes en el primario cuando se alimente una carga $R_L$ inductiva a trav\'es de un rectificador trif\'asico (carga no lineal trif\'asica balanceada).
        \item Se registrar\'an la tensi\'on del neutro en ausencia de conexi\'on y la corriente que lo atravesar\'a cuando est\'e conectado.
    \end{itemize}
    \item \textbf{Conexi\'on Y--Y alimentando una carga monof\'asica no lineal:}
    \begin{itemize}
        \item Se capturar\'a la forma de onda de las corrientes en el primario cuando se alimente una carga $R_L$ inductiva mediante un rectificador monof\'asico de onda completa.
    \end{itemize}
\end{enumerate}
\newpage

\section{Tabla de Resultados}

\renewcommand{\tablename}{Cuadro}
\setcounter{table}{0}

\begin{table}[H]
\centering
\small
\setlength{\tabcolsep}{4pt}
\caption{Mediciones en sistema de alimentaci\'on a tres hilos en vac\'io}
\begin{tabular}{|c|c|c|c|c|c|c|}
\hline
\multicolumn{4}{|c|}{Voltaje de fase (V)} & \multicolumn{3}{c|}{Voltaje de l\'inea (V)} \\ \hline
$V_a$ & $V_b$ & $V_c$ & $V_n$ & $V_{ab}$ & $V_{bc}$ & $V_{ac}$ \\ \hline
 & & & & & & \\ \hline
\end{tabular}
\end{table}

\begin{table}[H]
\centering
\small
\caption{Mediciones en sistema de alimentaci\'on a cuatro hilos en vac\'io}
\resizebox{\textwidth}{!}{%
\begin{tabular}{|c|c|c|c|c|c|c|c|c|}
\hline
\multicolumn{3}{|c|}{Voltaje de fase (V)} & \multicolumn{3}{c|}{Voltaje de l\'inea (V)} & \multicolumn{3}{c|}{Corrientes de fase (A)} \\ \hline
$V_a$ & $V_b$ & $V_c$ & $V_{ab}$ & $V_{bc}$ & $V_{ac}$ & $I_a$ & $I_b$ & $I_c$ \\ \hline
 & & & & & & & & \\ \hline
\end{tabular}%
}
\end{table}

\begin{table}[H]
\centering
\small
\caption{Mediciones en la conexi\'on Y--Y con carga trif\'asica no lineal con neutro}
\begin{tabular}{|c|c|c|}
\hline
$I_{\mathrm{neutro}}$ (A) & $I_{\mathrm{carga\ AC}}$ (A) & $I_{\mathrm{carga\ DC}}$ (A) \\ \hline
 & & \\ \hline
\end{tabular}
\end{table}

\begin{table}[H]
\centering
\small
\caption{Mediciones en la conexi\'on Y--Y con carga lineal trif\'asica balanceada en sistema de alimentaci\'on a cuatro hilos}
\resizebox{\textwidth}{!}{%
\begin{tabular}{|c|c|c|c|c|c|c|c|c|c|}
\hline
Carga (W) & $I_a$ (Alta) (A) & $I_b$ (Alta) (A) & $I_c$ (Alta) (A) & $V_a$ (Alta) (V) & $V_b$ (Alta) (V) & $V_c$ (Alta) (V) & $V_{L,a-b}$ (Alta) (V) & $V_{L,b-c}$ (Alta) (V) & $V_{L,a-c}$ (Alta) (V) \\ \hline
 & & & & & & & & & \\ \hline
 & & & & & & & & & \\ \hline
 & & & & & & & & & \\ \hline
 & & & & & & & & & \\ \hline
 & & & & & & & & & \\ \hline
 & & & & & & & & & \\ \hline
 & & & & & & & & & \\ \hline
 & & & & & & & & & \\ \hline
 & & & & & & & & & \\ \hline
 & & & & & & & & & \\ \hline
\end{tabular}%
}
\end{table}
\newpage

\section{Referencias}
\begin{enumerate}
    \item Nerio Ojeda, Juli\'an P\'erez. \textit{Pr\'actica 6. Transformador Monof\'asico. Rendimiento}. Universidad Central de Venezuela, Facultad de Ingenier\'ia, Escuela de Ingenier\'ia El\'ectrica, Departamento de Potencia, Laboratorio de M\'aquinas I, Revisi\'on 2012.
    \item Fraile Mora, Jes\'us. \textit{M\'aquinas El\'ectricas} (6.\textsuperscript{a} ed.). McGraw-Hill, 2008.
    \item Fitzgerald, A. E., Kingsley, C., Umans, S. D. \textit{M\'aquinas El\'ectricas}. McGraw-Hill, 6.\textsuperscript{a} ed., 2003.
    \item Chapman, S. J. \textit{M\'aquinas El\'ectricas}. McGraw-Hill, 5.\textsuperscript{a} ed., 2012.
\end{enumerate}

\end{document}
